% Based on Model 1 of "We Are a Learning Team!" by Urik Halliday

\model{Group vs Team}

Throughout the course, you will need to examine and process information, ask and answer questions, construct your own understanding, and develop new problem-solving skills.

\begin{center}
% source: http://www.clipartpanda.com/categories/group-prayer-images
\includegraphics[height=2.25in]{group1.png}
\hspace{0.5in}
% source: http://www.clipartpanda.com/categories/working-together-as-a-team
\includegraphics[height=2.25in]{group2.jpg}
\end{center}


\quest{8 min}


\Q What are some advantages to working in groups/teams?

\begin{answer}[5em]
You get to know other people and make new friends.
Different people have different backgrounds and skills.
The responsibility is shared.
\end{answer}


\Q What are some disadvantages to working in groups/teams?

\begin{answer}[5em]
Some group members may decide not to contribute.
One or two students may be absent.
People may not always get along with each other.
\end{answer}


\Q What is the difference between a group and a team?
Come up with a precise answer.

\begin{answer}[5em]
A group is students who just sit by each other and turn in the same assignment.
A team actually works together toward a common goal, drawing on each other's strengths.
\end{answer}


\Q How can working as a team help you accomplish the tasks described in \ref{group-vs-team.tex}?
Give at least two specific examples.

\begin{answer}[5em]
Working as a team makes it easier to examine and process information, because different people have different perspectives.
We can also develop new problem-solving skills by observing how each other approaches the problems.
\end{answer}
